(regular or otherwise). It has, on the other hand, the disadvantage (4.6) relative to @@M0@@, @@M1@@ that, as it stands, it is not covariant with respect to proper morphisms (even for smooth quasi-projective schemes over a field); it acquires this property only after tensoring with @@M2@@. Another advantage of @@M3@@ over its competitors @@M4@@ and @@M5@@ is the existence of a homomorphism from @@M6@@ to @@M7@@ that allows cohomological Chern classes to be interpreted as images of Chern classes with values in @@M8@@. In general this homomorphism does not factor through @@M9@@ (4.7), and it is at least doubtful that there is in general a homomorphism from @@M10@@ to @@M11@@ taking Chern classes to Chern classes (Cf. no.~5). Of course, these difficulties disappear after tensoring with @@M12@@, and one obtains a homomorphism
\sgaStableTarget{sga6.e14.eq.4-2}@@M13@@
taking Chern classes to Chern classes.

\sgaStableTarget{sga6.e14.para.4-2}\textbf{4.2.} To establish that the surjective homomorphism @@M14@@ of (4.1) is an isomorphism modulo torsion, one uses the Chern homomorphism (deduced from the theory of Chern classes [9])
@@M15@@
(where the second member is the Chern ring of @@M16@@ defined in Exp.~V), which is, as is easily seen, compatible with the topological filtration of @@M17@@ and the natural filtration of @@M18@@, and gives, on passing to the associated graded groups, a homomorphism of graded groups
\sgaStableTarget{sga6.e14.eq.4-3}@@M19@@
With this understood, we shall prove the formulas
\sgaStableTarget{sga6.e14.eq.4-4}@@M20@@
the first of which moreover implies the second, since @@M21@@ is surjective. These formulas clearly imply that @@M22@@ and @@M23@@ are isomorphisms modulo torsion and that, up to a factor @@M24@@ in degree @@M25@@, they are inverses of one another.

To prove the first formula (4.4), one can clearly suppose that @@M26@@ is defined by an irreducible closed subscheme @@M27@@ of @@M28@@ of codimension @@M29@@, in which case the formula is equivalent to
